\documentclass[11pt]{article}
\usepackage[margin=1in]{geometry}
\usepackage[T1]{fontenc}
\usepackage[utf8]{inputenc}
\usepackage{amsmath,amssymb,amsthm}
\usepackage{enumitem}

\title{ICPC World Finals 2021\\F. Islands from the Sky}
\author{}
\date{}

\begin{document}
\maketitle

\section*{Problem Summary}

Each flight travels along a straight segment in 3D. The camera always points vertically downward and has
aperture angle $\theta$. On the ground, that flight sees a strip around its projected path. An island is
successfully surveyed if \emph{one single flight} covers the whole polygon of that island.

We need the smallest possible angle $\theta$ for which every island is fully covered by at least one
flight, or report that this is impossible.

\section*{Shape Covered by One Flight}

Project a flight to the ground plane. Let its projected segment go from $A$ to $B$, and let the altitudes
at those endpoints be $z_A$ and $z_B$.

At any point along the flight, the camera sees points on the ground up to distance
\[
    z \tan \theta
\]
to each side of the projected path, measured perpendicularly. Because altitude changes linearly along the
flight segment, this half-width also changes linearly.

Therefore the visible ground region is a convex quadrilateral: a strip around $AB$ whose half-width is
$z_A \tan \theta$ at $A$ and $z_B \tan \theta$ at $B$.

\section*{A Vertex Formula}

Fix one flight and one island. Let
\begin{itemize}[leftmargin=*]
    \item $d$ be the unit direction vector of the projected flight;
    \item $n$ be a unit vector perpendicular to $d$.
\end{itemize}

For an island vertex $P$, write its coordinates relative to the flight start as
\[
    P-A.
\]
Then
\begin{itemize}[leftmargin=*]
    \item its coordinate \emph{along} the flight is
    \[
        s = (P-A)\cdot d;
    \]
    \item its perpendicular offset from the flight line is
    \[
        r = |(P-A)\cdot n|.
    \]
\end{itemize}

The point projects onto the flight segment if and only if
\[
    0 \le s \le |AB|.
\]
At that projected location, the altitude equals
\[
    z(s) = z_A + (z_B-z_A)\frac{s}{|AB|}.
\]
So this vertex is covered exactly when
\[
    r \le z(s)\tan\theta.
\]
Equivalently, this vertex requires
\[
    \tan\theta \ge \frac{r}{z(s)}.
\]

\section*{Why Checking Only Vertices Is Enough}

The visible region of one flight is convex. If every vertex of the island polygon lies inside that convex
region, then every island edge lies inside it as well, because each edge is a segment between two inside
points. Hence the whole island is covered.

So for a fixed island and flight, the minimum required value of $\tan\theta$ is just the maximum of
\[
    \frac{r}{z(s)}
\]
over all island vertices, provided every vertex projects onto the flight segment.

If even one vertex projects before $A$ or after $B$, then this flight can never cover the whole island,
regardless of the angle.

\section*{Algorithm}

\begin{enumerate}[leftmargin=*]
    \item For every island and every flight:
    \begin{itemize}[leftmargin=*]
        \item compute the projected flight direction and length;
        \item for each island vertex, compute its along-coordinate $s$ and perpendicular offset $r$;
        \item if some vertex has $s \notin [0,|AB|]$, mark this flight as impossible for that island;
        \item otherwise take the maximum required value
        \[
            \frac{r}{z(s)}.
        \]
    \end{itemize}
    \item For each island, keep the best flight, i.e. the minimum required $\tan\theta$.
    \item If some island has no feasible flight, output \texttt{impossible}.
    \item Otherwise the answer is
    \[
        \theta = \arctan\!\Bigl(\max_i \text{bestNeeded}_i\Bigr),
    \]
    converted to degrees.
\end{enumerate}

\section*{Correctness Proof}

We prove that the algorithm returns the correct answer.

\paragraph{Lemma 1.}
For a fixed flight, a ground point $P$ is covered if and only if its orthogonal projection onto the flight
path lies on the projected segment and its perpendicular distance from that segment is at most
$z(s)\tan\theta$, where $z(s)$ is the flight altitude above the projected point.

\paragraph{Proof.}
At every position along the flight, the camera covers exactly the interval of ground points perpendicular
to the flight path whose half-width is the current altitude times $\tan\theta$. Since altitude varies
linearly along the straight flight segment, the allowed perpendicular distance at projected position $s$ is
exactly $z(s)\tan\theta$. A point outside the endpoint slab cannot be seen because photographs are taken
only while the plane moves along the segment. \qed

\paragraph{Lemma 2.}
For a fixed island and flight, if every island vertex is covered, then the whole island is covered.

\paragraph{Proof.}
By Lemma 1 the visible region of one flight is a convex quadrilateral. Every edge of the island polygon is
a segment between two covered vertices, so the whole edge lies inside this convex region. Since the island
polygon is exactly the union of its edges and interior, the full island is covered. \qed

\paragraph{Lemma 3.}
For a fixed island and flight, if all island vertices project onto the flight segment, then the minimum
possible value of $\tan\theta$ that covers the island equals
\[
    \max_{\text{vertex }P}\frac{r(P)}{z(s(P))}.
\]

\paragraph{Proof.}
By Lemma 1, each vertex $P$ is covered exactly when
\[
    \tan\theta \ge \frac{r(P)}{z(s(P))}.
\]
All vertices are covered simultaneously exactly when $\tan\theta$ is at least the maximum of these lower
bounds. By Lemma 2 this is also sufficient for the entire island. \qed

\paragraph{Lemma 4.}
If some island vertex projects outside the flight segment, then that flight can never cover the whole
island.

\paragraph{Proof.}
By Lemma 1, any covered point must project onto the flight segment itself. A vertex projecting outside the
segment is therefore uncovered for every angle, so the whole island cannot be covered by that flight. \qed

\paragraph{Theorem.}
The algorithm outputs \texttt{impossible} exactly when no valid survey exists; otherwise it outputs the
smallest aperture angle that allows a complete survey.

\paragraph{Proof.}
By Lemmas 3 and 4, for each island-flight pair the algorithm computes exactly the minimum feasible
$\tan\theta$, or correctly marks the pair impossible. Choosing the minimum over flights gives the best way
to cover that island. Then the overall angle must be large enough to satisfy every island, so it is
determined by the maximum of these per-island requirements. Since $\arctan$ is strictly increasing, taking
the arctangent at the end gives exactly the smallest feasible angle. \qed

\section*{Complexity Analysis}

Let $V$ be the total number of island vertices. For each of the $m$ flights we inspect every island vertex
once, so the running time is $O(mV)$. With the given limits this is easily fast enough. The memory usage
is $O(V + m)$.

\section*{Implementation Notes}

\begin{itemize}[leftmargin=*]
    \item It is simpler to work with $\tan\theta$ during the whole computation and convert to degrees only
    once at the end.
    \item All geometry here is just dot products with one unit direction vector and one unit normal vector
    per flight.
\end{itemize}

\end{document}
